\section{低维结构为什么普遍存在}
\label{sec:12-Real-Analysis-Support/10-Essential-Structures/02}

手上有一份 \(64\times 64\) 的灰度图数据集，每张图是 \(\R^{4096}\) 里的一个点。按维数灾难的算术，要让最近邻检索在 \(4096\) 维里有意义，样本量得按指数增长，几万张图连零头都不够。可实践中最近邻检索在这几万张图上工作得相当好，检索回来的确实是相似的图。两件事只能有一个解释：这些向量没有填满 \(\R^{4096}\)，它们挤在一个薄得多的集合附近。问题是，凭什么？

\subsection{低维是一条经验命题，不是公理}

先把要讨论的对象说准。设数据分布 \(P\) 在 \(\R^d\) 上，若存在 \(k\) 维光滑子流形 \(\mathcal M\subset\R^d\) 与小的 \(\epsilon\)，使得 \(P\) 的绝大部分质量落在 \(\mathcal M\) 的 \(\epsilon\) 邻域内，且 \(k\ll d\)，就说这份数据满足流形假设。

这个陈述有三个可调的量：\(k\)、\(\epsilon\)、以及``绝大部分''是多少。三个量都不给，命题无法证伪；给了，它就是一条关于具体数据集的经验命题，可以在数据上检验，也可以在别的数据上不成立。把流形假设当作普遍规律来引用，是这一带最常见的滑坡。它在自然图像上有相当扎实的证据，在表格型业务数据上常常勉强，在纯计数数据上往往不成立。

\subsection{三类机制会把数据压到低维}

第一类是生成过程的自由度本来就少。把一个茶壶放在转台上，每转一度拍一张照，拍满 \(360\) 张。不管相机是百万像素还是千万像素，这批图像在像素空间里构成的集合是一条曲线：它由一个角度参数决定，内蕴维数是 \(1\)。写成 \(X=F(a)\)，\(a\in\mathcal A\subset\R^k\)，\(F\) 光滑且 Jacobian 处处满秩，则 \(F(\mathcal A)\) 是 \(k\) 维子流形——这一步是浸入定理，不是经验。经验的部分在于现实的 \(k\) 到底多大：转台只有一个自由度，而``自然场景照片''背后的自由度谁也说不清。

第二类是约束。上一节说对称性削减模型必须表达的自由度，同一句话对数据也成立：分子构象受键长与键角的强约束，可行构型集的维数远低于 \(3N\) 个坐标；刚体运动被正交条件 \(R^{\trans}R=I\) 限制在 \(SO(3)\) 上，\(9\) 个数只剩 \(3\) 个自由度。每一条硬约束都在环境空间里切掉一维。

第三类容易被漏掉：采样方式本身就是筛选。相机对准的是场景而非噪声，文本是人写给人看的，传感器读数是在设备正常工作时记录的。数据从来不是从 \(\R^d\) 的均匀分布里抽出来的，它是一个高度选择性的过程的输出。这条机制解释力很强，但也最不好量化——它说明低维结构与采集流程绑定，换一套采集流程，结论可能就变了。

\begin{center}
\begin{tikzpicture}[line width=0.8pt]
\node at (2.0,3.0) {\footnotesize 两个参数就能定位};
\foreach \v in {0,0.2,0.4,0.6,0.8,1}{
  \draw[gray,line width=0.4pt] plot[domain=0:1,samples=40]
    ({3*\x+1.0*\v},{1.5*\v+0.75*sin(\x*180)});
}
\foreach \u in {0,0.2,0.4,0.6,0.8,1}{
  \draw[gray,line width=0.4pt] ({3*\u},{0.75*sin(\u*180)})
    -- ({3*\u+1.0},{1.5+0.75*sin(\u*180)});
}
\foreach \u/\v in {0.08/0.15, 0.22/0.62, 0.35/0.28, 0.44/0.88, 0.57/0.45,
                   0.66/0.12, 0.73/0.71, 0.85/0.34, 0.93/0.80, 0.15/0.95,
                   0.50/0.05, 0.30/0.50}{
  \fill ({3*\u+1.0*\v},{1.5*\v+0.75*sin(\u*180)}) circle (1.7pt);
}
\begin{scope}[shift={(5.6,0)}]
\node at (1.6,3.0) {\footnotesize 两个参数定位不了};
\draw (0,0) -- (2.4,0) -- (2.4,2) -- (0,2) -- cycle;
\draw[gray] (0.85,0.7) -- (3.25,0.7) -- (3.25,2.7) -- (0.85,2.7) -- cycle;
\draw[gray] (0,0) -- (0.85,0.7);
\draw[gray] (2.4,0) -- (3.25,0.7);
\draw (2.4,2) -- (3.25,2.7);
\draw[gray] (0,2) -- (0.85,2.7);
\foreach \a/\b/\c in {0.35/0.40/0.15, 1.10/0.25/0.62, 1.85/0.55/0.30,
                      0.60/1.20/0.85, 1.45/1.05/0.10, 2.05/1.55/0.70,
                      0.25/1.70/0.40, 1.70/0.80/0.95, 0.95/1.60/0.22,
                      2.15/0.35/0.50, 0.50/0.90/0.55, 1.30/1.75/0.68,
                      1.95/1.25/0.05, 0.80/0.15/0.35, 2.25/1.85/0.90}{
  \fill ({\a+0.85*\c},{\b+0.7*\c}) circle (1.7pt);
}
\end{scope}
\end{tikzpicture}
\end{center}

\emph{左边的点云虽在三维中，却只沿一片二维曲面铺开；右边的点云填满整块立方体，任何二维坐标都留不住它。}

\subsection{谱衰减把``有多低维''变成可测量的}

线性版本最好办。把 \(n\) 个样本按行排成 \(X\in\R^{n\times d}\)，中心化后做奇异值分解，Eckart--Young 定理给出：秩 \(k\) 的最优逼近误差在 Frobenius 范数下等于 \(\big(\sum_{i>k}\sigma_i^2\big)^{1/2}\)。于是奇异值序列的衰减形状直接读出结论——衰减越陡，越少的方向就能兜住越多的能量。这条定理与它的用法在\hyperref[sec:03-Linear-Algebra/07-SVD-and-Data-Applications/03]{``低秩近似：用少数方向保留主要结构''}一节里有完整推导。所谓``数值秩''，就是把这条曲线在某个容差处截断得到的数。

麻烦在于反过来读。看到前几个特征值明显大于其余，很自然会认为找到了低维结构，但纯噪声也长这样。取 \(X\) 的元素为独立同分布、均值零方差 \(\sigma^2\) 的随机数，此时数据完全没有低维结构，可样本协方差的谱依然不是平的：当 \(d/n\to c\in(0,1)\)，Marchenko--Pastur 律说谱聚集在区间 \(\big[\sigma^2(1-\sqrt c)^2,\ \sigma^2(1+\sqrt c)^2\big]\) 上。\(d=100\)、\(n=400\) 时 \(c=0.25\)，区间是 \([0.25\sigma^2,\,2.25\sigma^2]\)：最大与最小特征值差九倍，而底下什么结构都没有。

这就把判据摆正了。判断低维结构存在与否，参照系是 MP 体的上边缘而不是``谱应该是平的''；只有明显冲出上边缘的那几个特征值才有资格被当作信号。\hyperref[ch:094]{``随机矩阵理论：高维协方差的谱失真''}一章把这套对照做成了可操作的流程，包括边缘处的涨落尺度和尖峰模型的相变阈值。

\subsection{反例与程度差别同样重要}

最干脆的反例是各向同性噪声。每个像素独立均匀取值生成的图像，内蕴维数就是 \(d\)，没有任何低维逼近能省下代价；上图右侧的立方体点云画的正是这种情形，它不是人为构造的病态例子。

更常见的是部分成立。股票日收益的协方差里，第一主成分（市场因子）远远冲出 MP 体，接下来几个行业因子勉强露头，剩下几百个特征值老老实实躺在体内。这份数据的诚实描述是``少数几维信号加上高维噪声''，而不是``低维流形''。做降维时保留前几维是对的，指望丢掉的部分是可忽略的残差就不对——残差的总方差往往占大头。

还有两处措辞需要收紧。其一，含噪数据的支撑集通常是全空间：\(X=F(A)+\varepsilon\) 中只要 \(\varepsilon\) 是满维的，严格意义下的低维子流形就不存在，低维只在``质量集中''这个意义上成立，而且成立与否依赖于分辨率——在 \(\epsilon\) 尺度以下看，数据是 \(d\) 维的。其二，不同类别的数据常有不同的内蕴维数，各分支还可能相交，此时``流形''这个词已经不合适，分层集合是更准确的说法。\hyperref[sec:14-Deep-Learning/10-Interpretability-Mathematics/03]{``高维数据的低维结构检验''}一节讨论了在这些情形下怎样做估计。

\subsection{哪些是定理，哪些是观察}

Eckart--Young 是定理，浸入定理是定理，Marchenko--Pastur 律在其独立同分布与四阶矩条件下也是定理。``自然数据集中在低维结构附近''是观察，逐数据集成立、逐数据集失败，强度差得很远。``深度网络的中间层是学到的充分统计量''则是启发：它给了一个有用的图像，但除非把任务和分布族都写死，否则没有可核对的陈述与之对应。

低维结构讲的是数据在空间里占多大一块，量的是集合的维数。下一节把镜头从数据移到模型：参数族本身可以看成一个几何对象，一旦给它配上合适的度量，统计与优化里原本分开处理的几件事会落到同一张图上。
